% Gerado por scripts/update_document_state.py; não editar manualmente.
\chapter*{Estado dos capítulos}
\addcontentsline{toc}{chapter}{Estado dos capítulos}
O quadro descreve o conteúdo desta versão. Capítulos planejados ainda não
integram o corpo da dissertação e não são apresentados como resultados obtidos.
\begingroup\small\singlespacing
\begin{longtable}{@{}p{0.08\textwidth}p{0.55\textwidth}p{0.20\textwidth}p{0.08\textwidth}@{}}
\toprule Cap. & Conteúdo & Estado & Marco \\
\midrule\endhead
1 & Introdução & Concluído & C02 \\ \addlinespace[5pt]
2 & Revisão bibliográfica & Concluído & C03 \\ \addlinespace[5pt]
3 & Fundamentos teóricos & Concluído & C04 \\ \addlinespace[5pt]
4 & Resposta de PTA & Concluído & C05 \\ \addlinespace[5pt]
5 & Metodologia de simulação & Concluído & C06 \\ \addlinespace[5pt]
6 & Inferência e validação & Concluído & C07 \\ \addlinespace[5pt]
7 & Resultados da compressão & Planejado & C08 \\ \addlinespace[5pt]
8 & Robustez estatística & Planejado & C09 \\ \addlinespace[5pt]
9 & Helicidade zero & Planejado & C10 \\ \addlinespace[5pt]
10 & Aplicação e alcance & Planejado & C11 \\ \addlinespace[5pt]
11 & Discussão & Planejado & C12 \\ \addlinespace[5pt]
12 & Conclusões & Planejado & C12 \\ \addlinespace[5pt]
\bottomrule\end{longtable}\endgroup
A introdução delimita as hipóteses e as condições de validação. A classificação
de um capítulo como concluído refere-se à entrega daquela etapa; revisões posteriores
serão incorporadas com rastreabilidade nos próximos releases.
